The radius $AB=8$ is perpendicular to tangent $BE$, so $BE=\sqrt{10^2-8^2}=6$. Since $BE:DE=2:3$, $BD=15$. Thus $AD=\sqrt{8^2+15^2}=17$. The diameter $XY$ is parallel to the chord of contact $BC$, hence perpendicular to $AD$. Therefore the triangle has base 16 and height 17, giving area $\boxed{136}$.
\begin{asy}
import graph;import geometry;usepackage("amsmath");size(7cm);
pair A=(0,0),D=(17,0),B=(64/17,120/17),C=(64/17,-120/17),E=B+.4*(D-B),X=(0,8),Y=(0,-8);draw(circle(A,8));draw(D--B^^D--C);draw(X--D--Y--cycle,blue);draw(A--E,dashed);label("$A$",A,(-1,0));label("$D$",D,(1,0));label("$B$",B,(0,1));label("$C$",C,(0,-1));label("$E$",E,(0,1));label("$X$",X,(0,1));label("$Y$",Y,(0,-1));
\end{asy}

